\documentclass[aspectratio=169,11pt]{beamer}
\usepackage[spanish,es-noquoting]{babel}
\usepackage[utf8]{inputenc}
\usepackage[T1]{fontenc}
\usepackage{lmodern,booktabs,tikz}
\usepackage{xurl}
\graphicspath{{../resultados/figuras/}{../resultados/}{./}}

% Colores Okabe-Ito (los mismos de las figuras)
\definecolor{azul}{HTML}{0072B2}
\definecolor{naranja}{HTML}{E69F00}
\definecolor{verde}{HTML}{009E73}
\definecolor{morado}{HTML}{CC79A7}
\definecolor{rojo}{HTML}{D55E00}
\definecolor{tinta}{HTML}{1F2933}

\usetheme{default}
\setbeamercolor{frametitle}{fg=azul}
\setbeamercolor{title}{fg=azul}
\setbeamercolor{structure}{fg=azul}
\setbeamercolor{normal text}{fg=tinta}
\setbeamerfont{frametitle}{series=\bfseries}
\setbeamertemplate{navigation symbols}{}
\setbeamertemplate{itemize items}[circle]
\setbeamertemplate{footline}{%
  \hfill{\scriptsize\color{gray}\insertshorttitle\quad\insertframenumber/\inserttotalframenumber}\hspace{1em}\vspace{0.6em}}

\newcommand{\repo}{\url{https://github.com/odvelez/StreetView}}
\newcommand{\quien}[1]{\begin{tikzpicture}[remember picture,overlay]\node[anchor=north east,font=\footnotesize,text=gray] at ([xshift=-1.2em,yshift=-0.9em]current page.north east) {#1};\end{tikzpicture}}
\newcommand{\cifra}[3][0.28\textwidth]{\begin{minipage}[t]{#1}\centering{\Huge\bfseries\color{azul}#2}\\[2pt]{\small #3}\end{minipage}}
\newcommand{\figura}[3]{\IfFileExists{../resultados/figuras/#1}{\includegraphics[#2]{#1}}{\fbox{\parbox{0.9\linewidth}{\centering\small #3}}}}

\title[Obras viales por etapas]{Obras viales por etapas}
\subtitle{Cómo planificar una obra que cierra calles\\Caso de estudio: canalización del arroyo El Salao, Soledad (Atlántico)}
\author{\small Oscar Vélez \and Gustavo de la Ossa \and Sebastián Blanco \and Karel Silva \and Santiago Lobo}
\institute{MAT 4021 Matemáticas Discretas, Universidad del Norte. Grupo 7}
\date{16 de noviembre de 2026}

\begin{document}

% ---------------------------------------------------------------- Oscar (0:00-2:00)
\begin{frame}
\titlepage
\vspace{-1em}
\begin{center}\small Repositorio: \repo\end{center}
\end{frame}

\begin{frame}{El problema}
\quien{Oscar}
\begin{columns}[T]
\column{0.55\textwidth}
\textbf{Por qué importa}
\begin{itemize}
\item El 4 de noviembre de 2026, una creciente de El Salao, en Soledad, dejó víctimas.
\item Los vecinos de ambos lados necesitan cruzarlo a diario.
\item Canalizarlo es una obra larga que cierra calles y puentes.
\end{itemize}
\column{0.42\textwidth}
\begin{block}{La pregunta}
Al planificar una obra que cierra calles, ¿cómo se divide en etapas y qué medidas temporales reducen más el desvío de quienes la rodean?
\end{block}
\end{columns}
\vfill
\small El Salao en el mapa: 2.386 m de obra, 5.859 esquinas y 8.827 tramos de calle.
\end{frame}

\begin{frame}{La idea: la obra se hace por etapas}
\quien{Oscar}
En cada etapa solo se cierran los cruces de ese tramo; el resto de la ciudad sigue abierta.
\medskip
\begin{columns}[T]
\column{0.24\textwidth}\textbf{\color{naranja}Paso peatonal provisional}\\\small Uno por cada 500 m de obra, mínimo uno por etapa.
\column{0.24\textwidth}\textbf{\color{morado}Cambio de sentido}\\\small Una calle de un solo sentido pasa a doble sentido.
\column{0.24\textwidth}\textbf{\color{azul}Ruta alterna}\\\small Para cada cruce cerrado, el camino más corto que queda. Se señaliza.
\column{0.24\textwidth}\textbf{\color{rojo}Control en esquinas}\\\small Las esquinas con más recorridos desviados.
\end{columns}
\vfill
\small Número de etapas = largo de la obra $\div$ 1.000 m, redondeado hacia arriba. No se proponen puentes vehiculares: costarían y tardarían casi lo que la obra.
\end{frame}

% ---------------------------------------------------------------- Karel (2:00-4:00)
\begin{frame}{Del mapa al grafo}
\quien{Karel}
\begin{columns}[T]
\column{0.55\textwidth}
\begin{itemize}
\item Esquina = nodo; cuadra = arista con su largo en metros (OpenStreetMap).
\item \textbf{Peatones:} grafo no dirigido. \textbf{Vehículos:} grafo dirigido (sentidos viales).
\item La obra \textbf{quita aristas}: las cuadras a menos de 10 m del eje, puentes existentes incluidos. Subir el peso sería inventar un número.
\item El lado de cada esquina: signo del producto cruz con el eje.
\end{itemize}
\column{0.42\textwidth}
\centering
\begin{tikzpicture}[xscale=1.1,yscale=1.4,esq/.style={circle,draw,fill=white,inner sep=2pt}]
\fill[gray!30] (-0.4,0.45) rectangle (4.4,0.55);
\draw (0,1) -- (4,1); \draw (0,0) -- (4,0);
\draw[rojo,dashed,very thick] (1,0) -- (1,1);
\draw[rojo,dashed,very thick] (3,0) -- (3,1);
\draw[naranja,line width=2pt] (2,0) -- (2,1);
\foreach \x in {0,1,2,3,4} {\node[esq] at (\x,1) {}; \node[esq] at (\x,0) {};}
\node[below,font=\scriptsize] at (2,-0.2) {\color{rojo}cerrada \quad \color{naranja}paso provisional};
\end{tikzpicture}
\end{columns}
\vfill
\[\Delta(a,b)=d_{\text{obra}}(a,b)-d_{\text{sin obra}}(a,b)\]
\end{frame}

\begin{frame}{Cómo se calcula}
\quien{Karel}
\begin{itemize}
\item \textbf{Dijkstra} (hecho a mano en \texttt{src/grafo.py}): camino mínimo sin obra y con obra.
\item \textbf{Fórmula de arco nuevo}: al agregar $x\to y$ de peso $w$,
\[d'(a,b)=\min\bigl(d(a,b),\ d(a,x)+w+d(y,b)\bigr)\]
Evalúa cada candidato sin repetir Dijkstra: de unos 36 s a 6--8 s, mismos resultados.
\item \textbf{Voraz}: en cada ronda, la medida que más baja el desvío medio. No garantiza el óptimo.
\item \textbf{Programación dinámica}: cortes cada 25 m; minimiza el mayor número de cruces en una etapa y luego la suma de cuadrados.
\end{itemize}
\end{frame}

\begin{frame}{Un caso resuelto a mano}
\quien{Karel}
\begin{columns}[c]
\column{0.45\textwidth}
\centering
\begin{tikzpicture}[xscale=1.1,yscale=1.6,esq/.style={circle,draw,fill=white,minimum size=6mm,inner sep=0pt}]
\fill[gray!30] (-0.4,0.45) rectangle (4.4,0.55);
\draw (0,1) -- (4,1); \draw (0,0) -- (4,0);
\draw[rojo,dashed,very thick] (0,0) -- (0,1);
\draw[rojo,dashed,very thick] (2,0) -- (2,1);
\draw[thick] (4,0) -- (4,1);
\node[esq] at (0,1) {$A$}; \node[esq] at (2,1) {$B$}; \node[esq] at (4,1) {$C$};
\node[esq] at (0,0) {$D$}; \node[esq] at (2,0) {$E$}; \node[esq] at (4,0) {$F$};
\node[above,font=\small] at (1,1.05) {100}; \node[above,font=\small] at (3,1.05) {100};
\node[below,font=\small] at (1,-0.05) {100}; \node[below,font=\small] at (3,-0.05) {100};
\node[left,rojo,font=\small] at (0,0.75) {80}; \node[left,rojo,font=\small] at (2,0.75) {50};
\node[right,font=\small] at (4,0.75) {50};
\end{tikzpicture}
\column{0.52\textwidth}
\small
La etapa cierra $A$--$D$ y $B$--$E$. Solo queda $C$--$F$.
\begin{itemize}
\item Desvío medio de los 9 pares: $970/9\approx\mathbf{107{,}8}$ m.
\item Paso en $B$--$E$: $d'(A,D)=\min(450,\,100+50+100)=250$; media \textbf{18,9 m}.
\item Paso en $A$--$D$: media \textbf{28,9 m}.
\end{itemize}
El voraz elige $B$--$E$. La prueba automática da lo mismo.
\end{columns}
\end{frame}

% ---------------------------------------------------------------- Gustavo (4:00-6:00)
\begin{frame}{El plan para El Salao: 2.386 m en tres etapas}
\quien{Gustavo}
\begin{columns}[c]
\column{0.58\textwidth}
\centering
\figura{plan_mapa.png}{width=\linewidth,height=0.68\textheight,keepaspectratio}{Mapa del plan (resultados/figuras/plan\_mapa.png)}
\column{0.40\textwidth}
\cifra[0.46\linewidth]{3}{etapas}\hfill\cifra[0.46\linewidth]{5}{pasos peatonales}\\[1em]
\cifra[0.46\linewidth]{3}{cambios de sentido}\hfill\cifra[0.46\linewidth]{18}{rutas alternas}\\[1em]
\small Etapas de 975, 850 y 561 m. 15 esquinas a controlar; 14 sin semáforo registrado.
\end{columns}
\end{frame}

\begin{frame}{A pie, el desvío baja en las tres etapas}
\quien{Gustavo}
\begin{columns}[c]
\column{0.55\textwidth}
\centering
\figura{desvio_a_pie.png}{width=\linewidth,height=0.7\textheight,keepaspectratio}{Barras de desvío a pie (resultados/figuras/desvio\_a\_pie.png)}
\column{0.43\textwidth}
\small
\begin{tabular}{@{}lrr@{}}
\toprule
Etapa & Con la obra & Con pasos\\
\midrule
1 (975 m) & 387 m & 48 m\\
2 (850 m) & 129 m & 28 m\\
3 (561 m) & 219 m & 56 m\\
\bottomrule
\end{tabular}
\medskip

Supuesto que más pesa: cada etapa cierra los puentes existentes. Si siguen abiertos, el desvío con la obra es 68, 56 y 36 m.
\end{columns}
\end{frame}

\begin{frame}{Rutas alternas y esquinas a controlar}
\quien{Gustavo}
\small
\begin{columns}[T]
\column{0.52\textwidth}
\textbf{Rutas alternas} (metros: sin obra $\to$ con la obra)\\[3pt]
\begin{tabular}{@{}llr@{}}
\toprule
Cruce cerrado & Etapa & Carro\\
\midrule
Calle 61 & 3 & 120 $\to$ 1.895\\
Vía local junto al eje & 2 & 102 $\to$ 1.776\\
Carrera 22A & 2 & 173 $\to$ 1.761\\
Calle 72 & 2 & 11 $\to$ 897\\
\bottomrule
\end{tabular}
\column{0.46\textwidth}
\textbf{Esquinas} (\% de recorridos desviados)\\[3pt]
\begin{tabular}{@{}lrl@{}}
\toprule
Esquina & \% & Semáforo\\
\midrule
Cl.~60 x Cra.~18 & 72,9 & no\\
Cl.~75 x Cra.~21 & 72,5 & no\\
Cl.~63 x Cra.~18 & 63,0 & sí\\
Cl.~73 x Cra.~19A & 61,1 & no\\
\bottomrule
\end{tabular}
\end{columns}
\vfill
Son candidatas: OpenStreetMap puede no tener todos los semáforos.
\end{frame}

% ---------------------------------------------------------------- Sebastián (6:00-8:00)
\begin{frame}{En carro, el cambio de sentido ayuda poco}
\quien{Sebastián}
\centering
\cifra{332 m}{etapa 1}\hspace{1em}\cifra{287 m}{etapa 2}\hspace{1em}\cifra{293 m}{etapa 3}\\[0.4em]
{\small desvío medio en carro con la obra}
\medskip
\begin{columns}[T]
\column{0.48\textwidth}
\begin{block}{Qué pasa}
Con el cambio de sentido queda en 332, 279 y 293 m. Su efecto real es dar salida: en la etapa 1, los pares sin camino pasan de 60 a 0.
\end{block}
\column{0.48\textwidth}
\begin{block}{Qué concluimos}
Con medidas temporales, el problema peatonal mejora mucho y el vehicular casi no. En carro se sostiene señalizar rutas y controlar esquinas.
\end{block}
\end{columns}
\end{frame}

\begin{frame}{¿Podemos confiar en los resultados?}
\quien{Sebastián}
\begin{itemize}
\item \textbf{Voraz contra fuerza bruta} ($k=2$): elige la misma pareja que el óptimo entre 1.225 parejas de pasos y 11.325 de sentidos. Un caso construido sí lo hace fallar.
\item \textbf{Sensibilidad}: con luz de 60 m o más, el primer paso es siempre el mismo y el desvío final queda entre 24 y 26 m.
\item \textbf{Control del lado}: el producto cruz coincide al 100~\% con las componentes conexas (1.229 esquinas).
\item \textbf{Pruebas automáticas}: \texttt{python -m unittest}, todas en verde.
\end{itemize}
\vfill
\small Las tres primeras, con toda la obra a la vez (\texttt{resultados/experimentos/}).
\end{frame}

\begin{frame}{Lo que este plan no mide}
\quien{Sebastián}
\begin{enumerate}
\item No hay datos de tránsito: no se mide congestión, solo desvíos en metros.
\item No se diseñan carriles ni tiempos de semáforo.
\item Se supone que cada etapa cierra los puentes existentes de su tramo.
\item Muestra de esquinas sobre OpenStreetMap, que puede estar incompleto.
\item El voraz no garantiza el óptimo.
\item Todo debe verificarse en terreno antes de proponerlo.
\end{enumerate}
\end{frame}

% ---------------------------------------------------------------- Santiago (8:00-10:00)
\begin{frame}{Demostración en vivo}
\quien{Santiago}
\begin{columns}[c]
\column{0.55\textwidth}
\begin{enumerate}
\item Cargar el ejemplo de El Salao.
\item Planificar y recorrer el plan paso a paso.
\item Abrir una ruta alterna: carro y a pie.
\item Otra obra: la calle 72 (ejemplo ilustrativo, no es una obra real).
\end{enumerate}
\column{0.42\textwidth}
\small
\begin{block}{Cómo correrlo}
\texttt{cd proyecto\_arroyos}\\
\texttt{pip install -r requirements.txt}\\
\texttt{python app.py}\\
\texttt{http://localhost:5000/?ejemplo}
\end{block}
\end{columns}
\end{frame}

\begin{frame}{Conclusiones}
\quien{Santiago}
\begin{enumerate}
\item Una obra que cierra calles se modela como quitar aristas; su efecto, como el desvío medido con Dijkstra.
\item A pie, los pasos provisionales funcionan: de 387 a 48 m en la etapa más difícil.
\item En carro, cambiar sentidos ayuda poco; se sostiene señalizar rutas alternas y controlar esquinas.
\item Dividir la obra en etapas limita el cierre a un tramo a la vez y sirve para cualquier obra.
\end{enumerate}
\vfill
\small Trabajo futuro: datos de tránsito, pesos en tiempo y muestra ponderada por población.
\end{frame}

\begin{frame}
\centering
{\Large\bfseries\color{azul} Gracias. ¿Preguntas?}\\[1.5em]
Código, datos, resultados y pruebas:\\[0.5em]
{\large \repo}\\[1.5em]
\small Datos del mapa: © colaboradores de OpenStreetMap (ODbL)
\end{frame}

\end{document}
